% 巨行星（综述）
% license CCBYSA3
% type Wiki

本文根据 CC-BY-SA 协议转载翻译自维基百科\href{https://en.wikipedia.org/wiki/Giant_planet}{相关文章}。

\begin{figure}[ht]
\centering
\includegraphics[width=8cm]{./figures/6d40f7553e942e4f.png}
\caption{太阳系的四颗巨行星如下：（上）木星与土星（气体巨行星）（下）天王星与海王星（冰巨行星）它们按照与太阳的距离顺序排列，并以真实色彩呈现。尺寸不按比例绘制。} \label{fig_Giant_1}
\end{figure}
巨行星（giant planet），有时也被称为类木行星（jovian planet）（Jove 是罗马神话中木星〔Jupiter〕的另一名称），是一类比地球大得多、类型多样的行星。巨行星通常主要由低沸点物质（挥发物）构成，而非岩石或其他固态物质，但“超级地球”（mega-Earths）这一类别也确实存在。在太阳系中共有四颗此类行星：木星、土星、天王星和海王星。此外，已有大量系外巨行星被发现。

巨行星有时被称为气体巨行星（gas giants），但许多天文学家现在将此称谓仅用于木星和土星，并将成分不同的天王星与海王星归类为冰巨行星（ice giants）。然而，这两种名称都可能导致误解：在太阳系中，所有巨行星主要由处于临界点之上的流体组成，在此状态下气相与液相不再具有明确区分。木星和土星主要由氢与氦构成，而天王星与海王星则由水、氨和甲烷构成。

关于低质量褐矮星与高质量气体巨行星（约 \(13\,M_{\mathrm{J}}\)）之间的定义差异仍存在争议。一种观点基于行星的形成过程；另一种观点则基于行星内部物理性质。争论的一部分涉及褐矮星是否必须在其演化史中的某一阶段经历过核聚变，这一点是否应被视为定义它们的必要条件。\(^\text{[1]}\)
\subsection{术语}
“气体巨行星”（gas giant）一词由科幻作家 James Blish 于 1952 年首次提出，最初用于指代所有巨行星。然而，这个名称在某种意义上并不准确，因为在这些行星的大部分体积中，压力高到足以使物质不再以气态存在。除大气的上层区域之外，行星内部的所有物质都可能处于临界点以上的状态，在该状态下液相与气相已无法区分。因此，“流体行星”（fluid planet）会是更精确的称呼。以木星为例，其中心附近存在金属氢，但其大部分体积仍由氢、氦以及微量其他气体构成，这些物质均处在其临界点以上。所有巨行星的可观测大气层（光学深度小于 1 的区域）相对行星半径而言极为薄弱，可能只延伸至中心约百分之一的位置。因此，可观测区域确实是气态的（与火星和地球不同，它们的大气层为观察其固体地壳提供了窗口）。

这一易引起误解的术语之所以流行，是因为行星科学家通常将 rock（岩质物质）、gas（气体）与 ice（冰质物质）作为行星成分的分类速记，与这些物质在行星内部的实际相态无关。在外太阳系中，氢与氦被归类为气体；水、甲烷与氨被归类为冰；硅酸盐与金属则被归类为岩石。若从行星深部内部的角度出发，“冰”基本指含氧与含碳物质，“岩石”指含硅物质，而“气体”则指氢和氦。由于天王星与海王星在许多方面与木星和土星显著不同，一些研究者开始仅将“气体巨行星”用于后两者。在此术语体系下，部分天文学家开始将天王星与海王星称为“冰巨行星”（ice giants），以强调其内部成分中冰（以流体形态存在）的占优势地位。

另一术语“类木行星”（jovian planet）源自罗马神话的木星神 Jupiter——其属格形式为 Jovis，故称 Jovian——意在表示这一类行星都与木星相似。

对质量足以点燃氘聚变的天体（在太阳成分条件下通常需超过 13 个木星质量）称为褐矮星（brown dwarf），其质量范围介于大型巨行星与最低质量恒星之间。13 个木星质量（\(13\,M_{\mathrm{J}}\)）的界限只是经验规则，而非具有精确物理意义的临界点。质量较大的天体会燃烧其大部分氘，而质量较小者只能燃烧极少量，13 \(M_{\mathrm{J}}\) 只是两者之间的大致分界。氘的燃烧量不仅取决于质量，也取决于行星的成分，特别是氦与氘的含量。Extrasolar Planets Encyclopaedia 将天体质量上限设置为 60 \(M_{\mathrm{J}}\)，而 Exoplanet Data Explorer 则采用 24 \(M_{\mathrm{J}}\) 的上限。\(^\text{[7][8]}\)

